% Options for packages loaded elsewhere
\PassOptionsToPackage{unicode}{hyperref}
\PassOptionsToPackage{hyphens}{url}
\documentclass[
  ignorenonframetext,
]{beamer}
\newif\ifbibliography
\usepackage{pgfpages}
\setbeamertemplate{caption}[numbered]
\setbeamertemplate{caption label separator}{: }
\setbeamercolor{caption name}{fg=normal text.fg}
\beamertemplatenavigationsymbolsempty
% remove section numbering
\setbeamertemplate{part page}{
  \centering
  \begin{beamercolorbox}[sep=16pt,center]{part title}
    \usebeamerfont{part title}\insertpart\par
  \end{beamercolorbox}
}
\setbeamertemplate{section page}{
  \centering
  \begin{beamercolorbox}[sep=12pt,center]{section title}
    \usebeamerfont{section title}\insertsection\par
  \end{beamercolorbox}
}
\setbeamertemplate{subsection page}{
  \centering
  \begin{beamercolorbox}[sep=8pt,center]{subsection title}
    \usebeamerfont{subsection title}\insertsubsection\par
  \end{beamercolorbox}
}
% Prevent slide breaks in the middle of a paragraph
\widowpenalties 1 10000
\raggedbottom
\AtBeginPart{
  \frame{\partpage}
}
\AtBeginSection{
  \ifbibliography
  \else
    \frame{\sectionpage}
  \fi
}
\AtBeginSubsection{
  \frame{\subsectionpage}
}
\usepackage{iftex}
\ifPDFTeX
  \usepackage[T1]{fontenc}
  \usepackage[utf8]{inputenc}
  \usepackage{textcomp} % provide euro and other symbols
\else % if luatex or xetex
  \usepackage{unicode-math} % this also loads fontspec
  \defaultfontfeatures{Scale=MatchLowercase}
  \defaultfontfeatures[\rmfamily]{Ligatures=TeX,Scale=1}
\fi
\usepackage{lmodern}
\ifPDFTeX\else
  % xetex/luatex font selection
\fi
% Use upquote if available, for straight quotes in verbatim environments
\IfFileExists{upquote.sty}{\usepackage{upquote}}{}
\IfFileExists{microtype.sty}{% use microtype if available
  \usepackage[]{microtype}
  \UseMicrotypeSet[protrusion]{basicmath} % disable protrusion for tt fonts
}{}
\makeatletter
\@ifundefined{KOMAClassName}{% if non-KOMA class
  \IfFileExists{parskip.sty}{%
    \usepackage{parskip}
  }{% else
    \setlength{\parindent}{0pt}
    \setlength{\parskip}{6pt plus 2pt minus 1pt}}
}{% if KOMA class
  \KOMAoptions{parskip=half}}
\makeatother
\usepackage{longtable,booktabs,array}
\usepackage{caption}
\captionsetup[table]{skip=6pt}
\newcounter{none} % for unnumbered tables
\usepackage{calc} % for calculating minipage widths
\usepackage{caption}
% Make caption package work with longtable
\makeatletter
\def\fnum@table{\tablename~\thetable}
\makeatother
\setlength{\emergencystretch}{3em} % prevent overfull lines
\providecommand{\tightlist}{%
  \setlength{\itemsep}{0pt}\setlength{\parskip}{0pt}}
% Auto-generated by infrastructure/rendering/_pdf_latex_helpers.py::extract_math_font_preamble.
% Minimal header passed to Pandoc via -H so Beamer slides pick up the same math font
% as the combined-PDF path. Adjust the manuscript's preamble.md to change the font globally.
\usepackage{unicode-math}
\setmathfont{latinmodern-math.otf}

% Slide layout defaults for warning-clean scientific decks.
\usepackage{etoolbox}
\IfFileExists{xurl.sty}{\usepackage{xurl}}{}
\IfFileExists{seqsplit.sty}{\usepackage{seqsplit}}{\newcommand{\seqsplit}[1]{#1}}
\protected\def\breakseq#1{\seqsplit{#1}}
\protected\def\breaktt#1{\begingroup\ttfamily\seqsplit{#1}\endgroup}
\setlength{\emergencystretch}{6em}
\tolerance=5000
\hbadness=10000
\hfuzz=1pt
\setlength{\tabcolsep}{2pt}
\AtBeginEnvironment{longtable}{\tiny\renewcommand{\arraystretch}{0.86}\setlength{\tabcolsep}{1pt}}
\AtBeginEnvironment{tabular}{\tiny\renewcommand{\arraystretch}{0.86}\setlength{\tabcolsep}{1pt}}
\AtBeginEnvironment{equation}{\tiny}
\AtBeginEnvironment{equation*}{\tiny}
\AtBeginEnvironment{align}{\tiny}
\AtBeginEnvironment{align*}{\tiny}
\AtBeginEnvironment{itemize}{\footnotesize}
\AtBeginEnvironment{enumerate}{\footnotesize}
\AtBeginEnvironment{description}{\footnotesize}
\setbeamerfont{caption}{size=\tiny}
\setbeamerfont{caption name}{size=\tiny}
\setbeamerfont{normal text}{size=\small}
\setbeamerfont{frametitle}{size=\small}
\setbeamerfont{section title}{size=\footnotesize}
\setbeamerfont{subsection title}{size=\footnotesize}
\setbeamertemplate{section page}{%
  \centering
  \begin{beamercolorbox}[sep=12pt,center,wd=\paperwidth]{section title}
    \parbox{0.86\paperwidth}{\centering\usebeamerfont{section title}\insertsection\par}
  \end{beamercolorbox}
}
\setbeamertemplate{subsection page}{%
  \centering
  \begin{beamercolorbox}[sep=8pt,center,wd=\paperwidth]{subsection title}
    \parbox{0.86\paperwidth}{\centering\usebeamerfont{subsection title}\insertsubsection\par}
  \end{beamercolorbox}
}
\setlength{\abovecaptionskip}{2pt}
\setlength{\belowcaptionskip}{0pt}

% Natbib and cross-reference fallbacks — slides don't load natbib
% or cleveref, but combined-PDF manuscript prose may emit these
% commands. The fallback renders citations as a bracketed key list
% and cross-references as detokenized labels so slides don't fail on
% undefined control sequences or raw underscores. \providecommand is
% a no-op if packages load later.
\providecommand{\citep}[1]{[#1]}
\providecommand{\citet}[1]{#1}
\providecommand{\citealp}[1]{#1}
\providecommand{\citeauthor}[1]{#1}
\providecommand{\citeyear}[1]{#1}
\providecommand{\cref}[1]{\texttt{\detokenize{#1}}}
\providecommand{\Cref}[1]{\texttt{\detokenize{#1}}}

\newtheorem{proposition}{Proposition}
\newtheorem{hypothesis}{Hypothesis}
\usepackage{bookmark}
\IfFileExists{xurl.sty}{\usepackage{xurl}}{} % add URL line breaks if available
\urlstyle{same}
% fallback for those not using the hyperref driver hyperxmp:
\makeatletter
\@ifundefined{xmpquote}{\newcommand{\xmpquote}[1]{#1}}{}
\makeatother
\hypersetup{
  hidelinks,
  pdfcreator={LaTeX via pandoc}}

\author{\texorpdfstring{}{}}
\date{}

\begin{document}

\section{Methods: Computational Kernels, Model Scope, and Quantitative
Provenance}\label{sec:methodology-kernels-provenance}

\subsection{Computational kernels as the sole authority for quantitative
examples}\label{computational-kernels-as-the-sole-authority-for-quantitative-examples}

\begin{frame}[fragile,allowframebreaks]{Computational kernels as the
sole authority for quantitative examples}
The \texttt{src/kernels} package implements the mathematics of Active
Inference as pure, numerically stable, fully typed functions in
natural-log (nats) units:

\end{frame}

\begin{frame}[fragile,allowframebreaks]{Computational kernels as the
sole authority for quantitative examples}
\begin{itemize}
\tightlist
\item
  \textbf{Information} -- surprise \(-\ln p\), Shannon entropy
  \(H(p)=-\sum_x p(x)\ln p(x)\), cross-entropy, Kullback-Leibler
  divergence \(\mathrm{KL}(p\|q)\), and mutual information, following
  the original and textbook information-theory treatments (Shannon 1948;
  Cover and Thomas 2006).
\item
  \textbf{Distributions} -- normalization, a precision-parameterized
  softmax, and Dirichlet expectations; these are the small probability
  utilities needed by the active-inference and variational examples
  (Beal 2003).
\item
  \textbf{Inference} -- exact discrete Bayesian posteriors and the
  variational free-energy identity
  \(F[q]=\mathbb{E}_{q(s)}[\ln q(s)-\ln p(o,s)]
  =D_{\mathrm{KL}}(q(s)\|p(s))-\mathbb{E}_{q(s)}[\ln p(o\mid s)]
  =D_{\mathrm{KL}}(q(s)\|p(s\mid o))-\ln p(o)\), together with the
  log-evidence and the evidence lower bound (Bishop 2006; Blei et al.
  2017; Buckley et al. 2017).
\item
  \textbf{Active inference} -- expected free energy \(G(\pi)\) with the
  convention that lower \(G\) is preferred, policy selection
  \(q(\pi)=\mathrm{softmax}(-\gamma G(\pi))\) where \(\gamma\) is policy
  precision, and a \breaktt{DiscreteActiveInferenceAgent} running a
  deterministic perception -\textgreater{} policy -\textgreater{} action
  loop (Friston et al. 2015; K. J. Friston et al. 2017; Da Costa et al.
  2020; Sajid et al. 2021; Millidge et al. 2021; Champion et al. 2026).
\end{itemize}
\end{frame}

\begin{frame}[allowframebreaks]{Scope of the implemented kernels}
\protect\phantomsection\label{scope-of-the-implemented-kernels}
The implemented kernels intentionally stop at a small, discrete,
one-step reference agent. Generalized free energy and action-oriented
model learning are important adjacent targets, but they remain scope
markers for future kernels rather than hidden features of this release
(Parr and Friston 2019; Tschantz et al. 2020).
\end{frame}

\begin{frame}[allowframebreaks]{Evidence boundary and analytic checks}
\protect\phantomsection\label{evidence-boundary-and-analytic-checks}
The resulting evidence is model-specific twice over. First, a tested
function establishes that the implementation matches its stated
mathematical identity under the declared assumptions and conventions.
Second, a fixed-input simulation shows how that selected model behaves
under its priors, likelihoods, preferences, precision, and policy set.
It does not show that the same assumptions describe a biological brain,
a human learner, or all formulations of Active Inference. This
distinction follows the separation between process-theory construction,
implementation choice, and empirical-status assessment in the
Active-Inference literature (K. Friston et al. 2017; Da Costa et al.
2020; Aguilera et al. 2022; Hodson et al. 2024).

\end{frame}

\begin{frame}[allowframebreaks]{Evidence boundary and analytic checks}
Kernel correctness is asserted against analytic identities: the entropy
of a uniform distribution over \(n\) outcomes equals \(\ln n\); KL is
non-negative and zero only at equality; the free-energy bound is tight
at the true posterior; and the expected-free- energy preference is
regime-dependent rather than unconditional. The reproducibility section
records the exact coverage and mock-free gates that must run before
treating this build as release evidence.
\end{frame}

\subsection{Statistical boundary: deterministic validation without
learner-outcome
inference}\label{statistical-boundary-deterministic-validation-without-learner-outcome-inference}

\begin{frame}[allowframebreaks]{Statistical boundary: deterministic
validation without learner-outcome inference}
Active Skillference does not report inferential statistics, learner
outcomes, effect sizes, \(p\)-values, confidence intervals, or
population-level empirical claims. Its evidence has three narrower
classes. First, deterministic enumeration validates graph shape,
source-matrix coverage, figure/artifact wiring, web math, roadmap state
and SkillTree export counts. Second, computational kernels and
simulation sweeps show model-specific mathematical behavior under fixed
inputs, such as noisy-sensor posterior sensitivity and the sample/commit
expected-free-energy boundary. Third, browser smoke checks and optional
online source-health checks are delivery and currentness evidence. Those
checks can reveal stale routes, missing artifacts, failed MathJax
rendering or unreachable source locators; they are not evidence that
learners improved or that a biological process theory is empirically
confirmed.
\end{frame}

\begin{frame}[allowframebreaks]{Statistical boundary: deterministic
validation without learner-outcome inference}

The next two layers apply that evidence discipline to authored content.
The first keeps numbers traceable to kernels; the second keeps citations
traceable to explicit support roles.
\end{frame}

\subsection{Provenance binding from kernel registry to learner-facing
prose}\label{provenance-binding-from-kernel-registry-to-learner-facing-prose}

\begin{frame}[fragile,allowframebreaks]{Provenance binding from kernel
registry to learner-facing prose}
Worked examples in \breaktt{src/kernels/examples.py} run the kernels on
fixed, documented inputs and register each result in a
\texttt{ClaimRegistry} under a stable name. Skill descriptions,
assessment stems, correct numeric answer options, and manuscript prose
reference these values only through \texttt{claim:NAME} tokens,
following the reproducible-reporting pattern of binding prose to
executable analysis rather than hand-maintained results (Xie 2015;
Sandve et al. 2013). Three gates enforce integrity:

\end{frame}

\begin{frame}[fragile,allowframebreaks]{Provenance binding from kernel
registry to learner-facing prose}
\begin{enumerate}
\tightlist
\item
  \textbf{Resolution} -- every token must resolve to a registered claim,
  or rendering raises.
\item
  \textbf{Fail-closed export} -- the SkillTree exporter calls the
  resolver; an unbacked token aborts the export with nothing written.
\item
  \textbf{Bare-number linter} -- any decimal result with two or more
  fractional digits sitting in plain prose (not inside a claim token or
  \texttt{\$...\$} math) is reported as a violation. Correct answer
  options are scanned as well, so a derived value cannot bypass
  provenance by hiding in an answer key; wrong distractors remain
  free-form.
\end{enumerate}

\end{frame}

\begin{frame}[fragile,allowframebreaks]{Provenance binding from kernel
registry to learner-facing prose}
We use ``provenance'' in a PROV-inspired sense: the system records
entities, activities, and agents sufficient for local traceability, but
it does not claim conformance to a full W3C PROV interchange
implementation. In Active Skillference terms, the entities are claim
records, rendered skills, manuscript fragments, figures, and SkillTree
exports; the activities are kernel execution, rendering, export, and
validation; and the agents are the source files and scripts that produce
each artifact. The PROV overview, data model, and constraints
recommendation supply the formal vocabulary for derivation, attribution,
bundles, and validity checks, while this project implements a
project-local claim registry and release ledger rather than a general
PROV service (Groth and Moreau 2013; Moreau and Missier 2013; Cheney et
al. 2013). Workflow Run RO-Crate and the related Provenance Run Crate
profile provide a useful interoperability target for the same boundary:
\end{frame}

\begin{frame}[fragile,allowframebreaks]{Provenance binding from kernel
registry to learner-facing prose}
they package workflow-run metadata, inputs, outputs, parameters, and
execution provenance as RO-Crate records. Active Skillference's
\texttt{claims.json}, manuscript variables, figure registry, and release
ledger play a narrower analogous role inside this repository. The
project does not emit a conforming RO-Crate, validate against those
profiles, or expose a general provenance interchange endpoint; the
comparison is a roadmap-level mapping from the local artifact chain to a
recognized workflow-provenance model ({Leo et al.} 2024).

\end{frame}

\begin{frame}[fragile,allowframebreaks]{Provenance binding from kernel
registry to learner-facing prose}
This mechanism sits in the lineage of literate and reproducible
computational reporting: Knuth made code and explanation a joint design
object, Sweave/knitr-style workflows bind prose to executable analysis,
and reproducible-research arguments treat rerunnable computational
claims as a minimum standard when independent replication is difficult
(Knuth 1984; Xie 2015; Peng 2011; Sandve et al. 2013). Active
Skillference moves that binding to a registry and build gate that
protects both the manuscript and the exported curriculum.

\end{frame}

\begin{frame}[fragile,allowframebreaks]{Provenance binding from kernel
registry to learner-facing prose}
Provenance binding and kernel testing are complementary halves of the
same pedagogical contract. Binding establishes that a value a learner
sees is a kernel's output rather than a hand-typed figure; the
deterministic kernel tests establish that the kernel computes the
intended quantity in the first place. A learner-facing number is
therefore only as sound as the kernel behind it, which is why the
integrity story rests on the tested kernels and not on the binding
alone: the binding makes a number traceable, and the tests are what make
that traceability worth anything. The export and render gates abort when
a kernel token is missing or a result was typed into prose, so an
unbacked number cannot reach a learner unnoticed.

\end{frame}

\begin{frame}[fragile,allowframebreaks]{Provenance binding from kernel
registry to learner-facing prose}
Numeric traceability is not enough for a scholarly curriculum. The prose
also needs to show which cited sources support lineage, method,
delivery, critique, and scope claims.
\end{frame}

\end{document}
